Large language models generate text token-by-token without any built-in mechanism for reflecting on what they have already written. Humans, by contrast, naturally pause between sentences to consider accuracy and coherence. We investigate whether prompting an LLM to generate explicit thinking tokens between every sentence---a technique we call \emph{sentence-level chain-of-thought} (\sentthink)---produces more careful text. Across three benchmarks using \gptfour, we find a striking trade-off: \sentthink significantly improves truthfulness on \tqa (4.69 vs.\ 4.50 out of 5.0, Wilcoxon $p{=}0.0007$, Cohen's $d{=}0.36$) but substantially reduces open-ended response quality on \mtbench (7.2--7.4 vs.\ 9.8--9.9 out of 10) because thinking tokens consume roughly half the output budget. On \gsm mathematical reasoning, \sentthink matches direct generation at 92\% accuracy. Contrary to our expectations, the resulting text is not more hedged or cautious---it is more concise and assertive, with fewer qualifiers and hedging words (all $p{>}0.3$). The interleaved thinking acts as a fact-checking filter, catching common misconceptions before they enter the response. Our results demonstrate that prompt-only interleaved deliberation is a viable zero-training technique for improving truthfulness in settings where response brevity is acceptable, while highlighting the fundamental token-budget constraint that separates prompting approaches from training-based latent reasoning methods.
